\sscp{The development of mid vowels in Wallacea}{13-01-05}
\citet[247--249]{Blust1993} lays out his arguments
and evidence for PMP with four vowels: *i,
*ə, *a, *u developing into PCEMP having six vowels with the addition
of *e and *o as: *i,
*e, *ə, *a, *o, *u.
He says these “\it{irregular} phonological changes [{\ldots}]
have considerable probative value for the CEMP hypothesis” (p. 246--247,
emphasis in the original).
The addition of the two new vowels is based on lexical developments in four words.
Blust’s arguments for mid-vowels in PCEMP provide a focused platform
for discussing the broader topic of mid vowels in Wallacean subgroups,
where they come from,
and whether there are patterns that might be useful for subgrouping.

First, we note that outside our target region,
\citet{Mills1975} reconstructed Proto-South Sulawesi as having
a six-vowel system including mid vowels with *i,
*e, *ɨ, *a, *o, *u.
We also note that Proto-Celebic is reconstructed as having mid
vowels \citep{Mead2003a,Mead2003b}, with PMP *-ay > *e,
and *-aw > *o,
and with Proto-Eastern Celebic further showing *ə > o,
and *-iq > *eq.
Given that mid-vowels are found just west of the Blust line,
development of mid-vowels would not qualify as
an exclusively shared innovation with subgrouping value. 

Before examining the specifics of Blust’s arguments
and evidence for PCEMP mid vowels,
it is worth noting some broad patterns in our target region relating
to common historical sources of \it{e}
and \it{o} in Austronesian languages in Wallacea.

\begin{itemize}
	\item	Most Austronesian languages in the region have a basic
				5-vowel system \it{(i, e, a, o, u}),
				with various conditioned allophones,
				and combinations or restrictions on various VV sequences.
				VV sequences of like vowels (\it{ii,
				ee, aa, oo, uu}) are common across the region
				and contrast both phonetically
				and phonemically with single V \citep{Grimes1991c,CulhaneEdwards2021,Tamelan2021}.
	\item	6-vowel systems (some with schwa \it{ə}) in a number of
				\tsc{\ili{Bima}-Lembata} languages are discussed in \srf{03-02-01} and \srf{03-03-01}.
	\item	Just outside the target region,
				several SHWNG languages have 7-vowel systems with \it{i,
				ɛ, e, a, ɔ, o, u}.
	\item	Some \tsc{East Bomberai} languages such as \ili{Irarutu} also
				appear to have 7-vowel systems (\citealp[651]{Voorhoeve1995a}; \citealp[43]{Jackson2014}).
				\tsc{East Bomberai} languages are in contact with SHWNG languages.
	\item	The \ili{Tarangan} languages in the \tsc{Aru} subgroup also have
				a 7-vowel system, with contrasting /e/,
				/ɛ/ and /o/, /ɔ/ (see \ref{ch:Aru}).
	\item Some languages \tsc{Southwest Maluku} also have multiple
				mid vowels, such as \ili{Emplawas} contrasting /e/ with /ɛ/.
				(But no /o/ vs. /ɔ/ contrast.)
	\item	\ili{As} shown in the following discussion,
				many languages in every subgroup in the region regularly get
				a phonemic /o/ vowel from an unconditioned split in PMP *u > \it{u, o}.
				Some get it from coalescence of PMP *-aw > \it{o},
				or from PMP *ə > \it{o}.
				A few also mirror *-aw > \it{o} with *au > \it{o,
				ou,} or \it{oo}, and with secondary vowel sequences from the
				loss of a historical consonant *a(C)u > \it{o,
				ou} or \it{oo}, (and a few > \it{oi}) or as a residual trace
				in the vowels from loss of PMP *w (e.g. *awa > \it{oa}).
				In the split of *u > \it{u,
				o}, which words in which languages get which vowels is often
				unpatterned and unpredictable (see examples later in this section),
				but it does mean that languages in different subgroups may surface
				with the same vowel through parallel but independent development.
				They get there via different paths.
				So the fact that a number of widely scattered languages in different
				Wallacean subgroups may collectively have \it{o} in the same
				word should not be surprising.
				An equal number of languages do not.
				The problem is, no consistent pattern has been found,
				so we are not dealing with regular sound correspondences in the rigorous scientific sense.
	\item	\ili{As} shown in \trf{13-03}
				and the following discussion,
				many languages in every subgroup in the region regularly get
				a phonemic /e/ vowel from PMP *ə > \it{e}.
				Some get it from coalescence of PMP *-ay > \it{e}.
				A few also mirror *-ay > \it{e} with *ai > \it{e,
				ei,} or \it{ee}, and secondary vowel sequences from the loss
				of a historical consonant *a(C)i > \it{e,
				ei} or \it{ee}, or as a residual trace in the vowels from loss
				of PMP *y (e.g.
				*aya > \it{ea, ae}).
				A few languages show a raising of *aq > \it{e}.
				The fact that a number of widely scattered languages can have
				\it{e} in the same word through independent parallel development
				should not be surprising.
				They got there via different paths.\footnote{
						For example,
						in \tsc{Sula-Buru} languages \ili{Buru}
						and \ili{Lisela}, PMP *ə > \it{e}.
						In \tsc{Taliabu} languages \ili{Taliabu}
						and \ili{Soboyo}, *ə > \it{o}, but *-əj > \it{o}, \it{e}.
						Thus there are some etyma,
						such as *pusəj ‘navel’,
						in which all four languages reflect final /e/,
						albeit through different historical processes.}
				Again,
				the problem is, no consistent pattern has been found,
				whether looking at the languages or the etyma.
	\item	The unconditioned split in PMP *u > \it{u,
				o} and *ə > \it{e} are separate phenomena.
				The list of languages that do the first only partially overlaps
				with the list of languages that do the second.
				For example in the \tsc{\ili{Ambon}-Seram} subgroup,
				there are some languages in which *u > \it{u, o}
				and *ə > \it{e}, others in which *u > \it{u, o}
				and *ə > \it{o}, and still others in which *u > \it{u, o}
				and *ə > \it{o, e}.
	\item	There are several cases of high vowels (*i,
				*u) lowering (to \it{e, o}) before PMP *R.
				This is found in parts of \tsc{\ili{Bima}-Lembata} (\trf{03-28}),
				\tsc{Rote-Meto} (\trf{04-04}; \citealp[84f]{Edwards2021}),
				and \tsc{\ili{Central} Timor} (\trf{05-11}).
	\item	In a few cases vowels can be raised or lowered due to open
				or closed (historical) syllables,
				harmony of vowel height,
				or disharmony of vowel height (see \trf{09-05} and surrounding discussion).
\end{itemize}

Looking at the collective evidence that Blust provides for CEMP
with specific attention to the support for mid vowels,
he concludes \citeyear[740]{Blust2013}, “Given the high quality of evidence
for CEMP it probably is safe to say that this subgroup is as
well-established as any large AN subgroup apart from Oceanic.”
Yet there are a number of discrepancies that undermine how confidently
we can accept that statement.